% !TeX program = lualatex
% !TeX encoding = utf8
% !TeX spellcheck = uk_UA
% !TeX root = ../EMConspect.tex

%=========================================================
\chapter{Енергія електричного поля}\label{\currfilebase}
\graphicspath{{\currfilebase/Pictures}}
%=========================================================
\epigraph{Енергія електричного поля --- це не властивість зарядів, а властивість простору навколо них: поле не лише переносить взаємодію між
	зарядами, а й саме є носієм і сховищем енергії.
}{Переосмислення ідей сучасної фізики}


%% --------------------------------------------------------
\section{Електрична ємність}
%% --------------------------------------------------------

\subsection*{Ємність відокремленого провідника}

Нехай усамітнений (далекий від інших тіл) провідник несе заряд $q$; тоді, як ми знаємо, весь цей заряд розподіляється по його поверхні, а сам
провідник --- еквіпотенціальне тіло з деяким потенціалом $\phi$. Якщо заряд збільшити в $n$ разів, то, за принципом суперпозиції, поле в кожній точці
простору також зросте в $n$ разів, а отже, в $n$ разів зросте й робота з переміщення пробного заряду з поверхні провідника на нескінченність, тобто й
сам потенціал $\phi$. Це означає, що відношення
\begin{equation}\label{eq:CapacitorDef}
	\tcbhighmath{
		C = \frac{q}{\phi}
	}
\end{equation}
не залежить від заряду провідника і визначається лише його формою, розмірами та властивостями середовища, що його оточує. Цю величину називають
\emphx{електричною ємністю провідника}.

\begin{Attention}
	Ємність провідника --- суто геометрична характеристика: вона залежить від форми й розмірів тіла, а також від діелектричної проникності
	середовища, але \emphx{не залежить} ні від заряду, ні від матеріалу самого провідника. Означення~\eqref{eq:CapacitorDef} коректне лише для
	\emphx{усамітненого} провідника, тобто коли поблизу немає інших заряджених тіл, здатних істотно змінити його потенціал.
\end{Attention}

\subsection*{Взаємна ємність двох провідників. Конденсатор}

Якщо поблизу є інші провідники, то потенціал даного тіла залежить уже не тільки від його власного заряду, а й від зарядів і розташування сусідніх
тіл, тому означення~\eqref{eq:CapacitorDef} втрачає сенс. Розглянемо систему з \emphx{двох} провідників, на яких розміщено рівні за величиною та
протилежні за знаком заряди $+q$ і $-q$. Різниця потенціалів $\phi_+ - \phi_-$ між провідниками, за тим самим міркуванням про суперпозицію, як і
вище, пропорційна заряду $q$. Відношення
\begin{equation}\label{eq:MutualCapacitanceDef}
	\tcbhighmath{
		C = \frac{q}{\phi_+ - \phi_-}
	}
\end{equation}
називають \emphx{взаємною ємністю} пари провідників.

\begin{Attention}
	\emphx{Конденсатор} --- це система з двох провідників (\emphx{обкладок}), розташованих настільки близько один до одного (і, як правило,
	розділених тонким шаром діелектрика), що практично все електричне поле зосереджене у вузькому проміжку між обкладками. Завдяки такому
	\emphx{екрануванню} поля взаємна ємність конденсатора, на відміну від ємності одинокого провідника, майже не залежить від наявності інших тіл
	поблизу --- саме тому вона є зручною і відтворюваною характеристикою.
\end{Attention}

\subsection*{З'єднання конденсаторів}

На практиці конденсатори часто з'єднують у батареї --- послідовно або паралельно, --- щоб отримати потрібну сумарну ємність.

\emphx{Послідовне з'єднання.} Нехай $n$ конденсаторів з'єднано послідовно (рис.~\ref{pic:CapacitorsSeries}). Оскільки обкладки сусідніх
конденсаторів, з'єднані між собою провідником, електрично ізольовані від решти кола, повний заряд кожної внутрішньої пари обкладок мусить
залишатися нульовим: якщо на крайню ліву обкладку батареї нанести заряд $+q$, то шляхом електростатичної індукції той самий за модулем заряд $q$
з'явиться на \emphx{кожному} конденсаторі батареї. Загальна різниця потенціалів (напруга) на батареї дорівнює сумі напруг на окремих конденсаторах:
\begin{equation*}
	U = U_1 + U_2 + \ldots + U_n = \frac{q}{C_1} + \frac{q}{C_2} + \ldots + \frac{q}{C_n}.
\end{equation*}
Оскільки ємність батареї за означенням $C = q/U$, маємо
\begin{equation}\label{eq:SeriesCap}
	\tcbhighmath{
		\frac1C = \sum\limits_{i=1}^n \frac1{C_i}.
	}
\end{equation}

\begin{figure}[h!]\centering
%	\localinput{capacitors-series}
	\caption{Послідовне з'єднання конденсаторів: заряд на всіх обкладках однаковий, напруги
		додаються\label{pic:CapacitorsSeries}}
\end{figure}

\emphx{Паралельне з'єднання.} Нехай тепер $n$ конденсаторів з'єднано паралельно (рис.~\ref{pic:CapacitorsParallel}): усі ліві обкладки з'єднані між
собою і мають спільний потенціал, так само як і всі праві. Тому напруга на кожному конденсаторі однакова і дорівнює загальній напрузі батареї $U$, а
заряди на окремих конденсаторах додаються:
\begin{equation*}
	q = q_1 + q_2 + \ldots + q_n = C_1 U + C_2 U + \ldots + C_n U.
\end{equation*}
Звідси
\begin{equation}\label{eq:ParallelCap}
	\tcbhighmath{
		C = \sum\limits_{i=1}^n C_i.
	}
\end{equation}

\begin{figure}[h!]\centering
%	\localinput{capacitors-parallel}
	\caption{Паралельне з'єднання конденсаторів: напруга на всіх конденсаторах однакова, заряди
		додаються\label{pic:CapacitorsParallel}}
\end{figure}

\begin{Attention}
	Послідовне з'єднання завжди \emphx{зменшує} сумарну ємність порівняно з найменшою з ємностей $C_i$ (аналогія з опорами резисторів у
	паралельному з'єднанні), тоді як паралельне з'єднання завжди \emphx{збільшує} сумарну ємність. Це протилежно до того, як поводяться опори при
	послідовному й паралельному з'єднанні резисторів, і пов'язано з тим, що ємність --- це <<здатність накопичувати заряд>>, а не <<перешкода струму>>.
\end{Attention}

\subsection*{Обчислення ємності найпростіших конденсаторів}

Розглянемо кілька типових задач на обчислення ємності, користуючись означеннями~\eqref{eq:CapacitorDef} і~\eqref{eq:MutualCapacitanceDef}, а також
результатами про поле в діелектрику (гл.~\ref{FileldInDielectrics}).

\emphx{Ємність усамітненої кулі.} Провідна куля радіусом $R$, що несе заряд $Q$, створює поза собою поле точкового заряду, а її власний потенціал
дорівнює потенціалу на поверхні, $\phi = Q/R$. Звідси відразу
\begin{equation}\label{eq:SphereCap}
	\tcbhighmath{
		C = R.
	}
\end{equation}
Цей простий результат корисно пам'ятати як орієнтир: ємність провідного тіла має порядок величини його лінійних розмірів.

\emphx{Плоский конденсатор.} Нехай обкладки --- дві паралельні пластини площею $S$ кожна, розташовані на відстані $d \ll \sqrt{S}$ одна від одної
(щоб знехтувати крайовими ефектами), а простір між ними заповнений діелектриком з проникністю $\epsilon$. Поле між обкладками однорідне,
$E = 4\pi\sigma/\epsilon = 4\pi q /(\epsilon S)$, тому різниця потенціалів $\Delta\phi = Ed = 4\pi qd/(\epsilon S)$, і
\begin{equation}\label{eq:FlatCap}
	\tcbhighmath{
		C = \frac{\epsilon S}{4\pi d}.
	}
\end{equation}

\emphx{Сферичний конденсатор.} Обкладки --- дві концентричні сфери радіусів $R_1 < R_2$, простір між якими заповнений діелектриком з проникністю
$\epsilon$. За теоремою Гаусса поле між обкладками таке саме, як і поле точкового заряду $q$ у діелектрику, $E = q/(\epsilon r^2)$. Різниця
потенціалів
\begin{equation*}
	\Delta\phi = \int\limits_{R_1}^{R_2} E\, dr = \frac{q}{\epsilon}\left(\frac1{R_1} - \frac1{R_2}\right) = \frac{q}{\epsilon}\frac{R_2 - R_1}{R_1
		R_2},
\end{equation*}
звідки
\begin{equation}\label{eq:SphericalCap}
	\tcbhighmath{
		C = \frac{\epsilon R_1 R_2}{R_2 - R_1}.
	}
\end{equation}
Легко переконатись, що при $R_2 \to \infty$ (друга обкладка йде на нескінченність) формула~\eqref{eq:SphericalCap} переходить у
формулу~\eqref{eq:SphereCap} для усамітненої кулі (з точністю до множника $\epsilon$), як і повинно бути.

\emphx{Циліндричний конденсатор.} Обкладки --- дві коаксіальні циліндричні поверхні радіусів $R_1 < R_2$ і довжини $\ell \gg R_2$ (щоб знехтувати
крайовими ефектами), простір між ними заповнений діелектриком з проникністю $\epsilon$. За теоремою Гаусса поле між обкладками
$E = \dfrac{2q}{\epsilon \ell r}$ (поле нескінченної зарядженої нитки в діелектрику), тому
\begin{equation*}
	\Delta\phi = \int\limits_{R_1}^{R_2} E\, dr = \frac{2q}{\epsilon\ell}\ln\frac{R_2}{R_1},
\end{equation*}
звідки
\begin{equation}\label{eq:CylindricalCap}
	\tcbhighmath{
		C = \frac{\epsilon \ell}{2\ln\left(R_2/R_1\right)}.
	}
\end{equation}

\subsection*{Взаємна ємність двох віддалених провідників}

Розглянемо систему з двох металевих куль, що мають власні ємності $C_1$ і $C_2$ (в сенсі означення~\eqref{eq:CapacitorDef}, тобто ємності кожної
кулі, якби вона була усамітненою) і перебувають на відстані $b$ одна від одної, причому $b \gg C_1, C_2$ (тобто $b$ набагато перевищує лінійні
розміри обох куль, так що одну відносно іншої можна вважати точковою, рис.~\ref{pic:TwoSpheres}).

\begin{SCfigure}[1.5][h!]\centering
	\localinput{two-spheres.tikz}
	\caption{Дві далекі одна від одної заряджені кулі\label{pic:TwoSpheres}}
\end{SCfigure}

Нанесемо заряд $+q$ на першу кулю і $-q$ --- на другу. Потенціал кожної кулі складається з потенціалу від власного заряду і від (малого) впливу
заряду сусідньої кулі, яку на такій відстані можна вважати точковим зарядом:
\begin{equation*}
	\begin{cases}
		\phi_+ = \dfrac{q}{C_1} - \dfrac{q}{b}, \\[4pt]
		\phi_- = -\dfrac{q}{C_2} + \dfrac{q}{b}.
	\end{cases}
\end{equation*}
Різниця потенціалів $\phi_+ - \phi_- = q\left(\dfrac1{C_1} + \dfrac1{C_2} - \dfrac2b\right)$, тому за
означенням~\eqref{eq:MutualCapacitanceDef} взаємна ємність системи дорівнює
\begin{equation*}
	C = \frac{q}{\phi_+ - \phi_-} = \frac{C_1 C_2}{C_1 + C_2 - \dfrac{2C_1C_2}{b}}.
\end{equation*}
Розкладаючи цей вираз у ряд за малим параметром $C_1C_2/[(C_1+C_2)b] \ll 1$, отримуємо наближену формулу
\begin{equation}\label{eq:TwoSpheresCap}
	\tcbhighmath{
		C \approx \frac{C_1C_2}{C_1+C_2}\left(1 + \frac{C_1C_2}{C_1+C_2}\frac{2}{b}\right).
	}
\end{equation}
У головному наближенні ($b \to \infty$) взаємна ємність двох далеких провідників дорівнює ємності їхнього послідовного з'єднання, а наступний доданок
описує невелику поправку, пов'язану з взаємним впливом куль одна на одну.

%% --------------------------------------------------------
\section{Енергія електричного поля}
%% --------------------------------------------------------

\subsection*{Взаємна енергія системи точкових зарядів}

Щоб зблизити два точкових заряди $q_1$ і $q_2$ з нескінченності до відстані $r_{12}$ один від одного, потрібно виконати проти сил поля роботу
$A = q_1q_2/r_{12}$. Цю роботу і ототожнюють з \emphx{взаємною енергією} пари зарядів: $U_{12} = q_1q_2/r_{12}$. Для системи з $n$ точкових зарядів
(рис.~\ref{pic:ChargesTriangle}) сумарна взаємна енергія набирається з енергій усіх можливих пар:
\begin{equation}\label{eq:PairEnergy}
	\tcbhighmath{
		U = \frac12 \sum\limits_{i=1}^n \sum\limits_{\substack{j=1\\ j\neq i}}^n \frac{q_iq_j}{r_{ij}},
	}
\end{equation}
де множник $1/2$ компенсує подвійний перелік кожної пари (доданок з індексами $i,j$ і доданок з індексами $j,i$ описують одну й ту саму пару).

\begin{SCfigure}[1.5][h!]\centering
	\localinput{charges-triangle.tikz}
	\caption{Система трьох точкових зарядів: взаємна енергія набирається із внесків усіх пар\label{pic:ChargesTriangle}}
\end{SCfigure}

Формулу~\eqref{eq:PairEnergy} зручно переписати через потенціал. Потенціал поля в точці розташування $i$-го заряду, створюваний \emphx{усіма іншими}
зарядами (крім самого $i$-го, поле якого в точці свого ж розташування не визначене),
\begin{equation*}
	\phi_i = \sum\limits_{\substack{j=1 \\ j \neq i}}^n \frac{q_j}{r_{ij}},
\end{equation*}
тому
\begin{equation}\label{eq:SystemEnergyViaPotential}
	\tcbhighmath{
		W = \frac12 \sum\limits_i q_i \phi_i.
	}
\end{equation}

\subsection*{Електростатична енергія тіла з неперервним розподілом заряду}

Для тіл з неперервним розподілом заряду --- об'ємною густиною $\rho(\vect r)$ та (якщо є) поверхневою густиною $\sigma(\vect r)$ ---
формулу~\eqref{eq:SystemEnergyViaPotential} узагальнюють природним чином, замінюючи суму на інтеграл, а $q_i$ --- на елемент заряду $\rho\, dV$ або
$\sigma\, dS$:
%(рис.~\ref{pic:BodyChargeEnergy})
\begin{equation}\label{eq:BodyEnergy}
	\tcbhighmath{
		W = \frac12 \iiint\limits_V \rho(\vect r)\phi(\vect r)\, dV + \frac12 \iint\limits_S \sigma(\vect r)\phi(\vect r)\, dS.
	}
\end{equation}

%\begin{SCfigure}[1][h!]\centering
%	\localinput{body-charge-energy.tikz}
%	\caption{Тіло з неперервно розподіленим зарядом\label{pic:BodyChargeEnergy}}
%\end{SCfigure}

\emphx{Приклад 1.} Нехай куля радіуса $R$ несе повний заряд $Q$, рівномірно розподілений по \emphx{поверхні}. Потенціал усюди на поверхні однаковий і
дорівнює $\phi = Q/R$, тому
\begin{equation}\label{eq:SphereEnergySurface}
	W = \frac12 \iint\limits_S \sigma\phi\, dS = \frac12\, \sigma\, \frac{Q}{R}\cdot 4\pi R^2 = \frac{Q^2}{2R}.
\end{equation}

\emphx{Приклад 2.} Нехай тепер той самий заряд $Q$ розподілений рівномірно по \emphx{об'єму} кулі. Обчислення (з урахуванням того, що потенціал
всередині рівномірно зарядженої кулі залежить від $r$) дає інший результат:
\begin{equation}\label{eq:SphereEnergyVolume}
	W = \frac35\, \frac{Q^2}{R}.
\end{equation}

\begin{Attention}
	Той самий сумарний заряд $Q$ на тілі того самого радіуса $R$ дає \emphx{різну} енергію залежно від того, як саме заряд розподілений усередині
	тіла (формули~\eqref{eq:SphereEnergySurface} і~\eqref{eq:SphereEnergyVolume} відрізняються множником $6/5$). Це наштовхує на природне запитання:
	чи можна вважати енергію \emphx{властивістю самих зарядів}, чи, може, правильніше говорити про енергію, \emphx{розподілену в просторі} --- там,
	де існує поле? Відповідь на це запитання дає перехід до об'ємної густини енергії поля, розглянутий далі.
\end{Attention}

\subsection*{Енергія зарядженого конденсатора}

Розглянемо процес зарядки конденсатора як послідовне перенесення малих порцій заряду $dq$ з однієї обкладки на іншу (рис.~\ref{pic:CapacitorDq}).
Коли на обкладках уже є заряд $q$ (а різниця потенціалів дорівнює $\Delta\phi = q/C$), перенесення чергової порції $dq$ вимагає роботи
$dW = \Delta\phi\, dq = \dfrac{q\,dq}{C}$. Інтегруючи від нульового заряду до кінцевого значення $q$, знаходимо енергію зарядженого конденсатора:
\begin{equation}\label{eq:CapacitorEnergy}
	\tcbhighmath{
		W = \frac{q^2}{2C} = \frac{C\Delta\phi^2}{2}.
	}
\end{equation}

\begin{SCfigure}[1.5][h!]\centering
	\localinput{capacitor-dq.tikz}
	\caption{Перенесення заряду $dq$ під час зарядки конденсатора\label{pic:CapacitorDq}}
\end{SCfigure}

Застосуємо цей результат до плоского конденсатора. Оскільки поле в ньому однорідне, $\Delta\phi = Ed$, а ємність дається
формулою~\eqref{eq:FlatCap}, маємо
\begin{equation*}
	W = \frac12 C \Delta\phi^2 = \frac12\,\frac{\epsilon S}{4\pi d}\,(Ed)^2 = \frac{\epsilon E^2}{8\pi}\, V,
\end{equation*}
де $V = Sd$ --- об'єм простору між обкладками, тобто об'єм, що його займає поле. Звідси природно ввести \emphx{об'ємну густину енергії поля}
\begin{equation}\label{eq:CapacitorEnergyDensity}
	\tcbhighmath{
		w = \frac{\epsilon E^2}{8\pi} = \frac{ED}{8\pi}
	}
\end{equation}
(останню рівність отримано підстановкою $D = \epsilon E$, див.~формулу~\eqref{eq:DepsE} з попереднього розділу).

\begin{Attention}
	Формула~\eqref{eq:CapacitorEnergyDensity} виражає енергію конденсатора \emphx{не} через заряди на обкладках, а через характеристики поля $\Efield$ і
	$\Dfield$ у тій області простору, де це поле існує. Це дає підставу для важливої фізичної інтерпретації: саме \emphx{електричне поле} є носієм
	взаємодії між зарядами і, водночас, носієм енергії --- енергія \emphx{локалізована} в тих областях простору, де є поле (у плоскому конденсаторі
	--- виключно в проміжку між обкладками), а не \enquote{приписана} безпосередньо зарядам.
\end{Attention}

\subsection*{Об'ємна густина енергії електричного поля (загальний випадок)}\label{sec:FieldEnergyDensity}

Формулу~\eqref{eq:CapacitorEnergyDensity} отримано для окремого випадку однорідного поля в конденсаторі. Покажемо, що вона має набагато загальніший
характер і справедлива для \emphx{будь-якого} електростатичного поля, однорідного чи ні.

Нехай заряд системи змінюється на нескінченно малу величину $\delta\rho(\vect r)$; відповідна зміна енергії, згідно з~\eqref{eq:BodyEnergy},
дорівнює
\begin{equation*}
	\delta W = \frac12 \iiint\limits_V \delta\rho(\vect r)\, \phi(\vect r)\, dV.
\end{equation*}
За теоремою Гаусса для вектора $\Dfield$ у диференціальній формі, $\delta\rho = \dfrac1{4\pi}\divg\delta\Dfield$, тому
\begin{equation*}
	\delta W = \frac1{8\pi} \iiint\limits_V \left(\divg \delta\Dfield\right) \phi\, dV.
\end{equation*}
Скористаємось векторною тотожністю $\divg(\phi\, \delta\Dfield) = \phi\, \divg\delta\Dfield + \delta\Dfield\cdot\vect\nabla\phi$, звідки
\begin{equation*}
	\delta W = \frac1{8\pi} \iiint\limits_V \divg\left(\phi\, \delta\Dfield\right) dV - \frac1{8\pi} \iiint\limits_V \delta\Dfield\cdot\vect\nabla\phi\, dV.
\end{equation*}
Поширимо інтегрування на весь (нескінченний) простір. Перший доданок, за теоремою Остроградського--Гаусса, дорівнює потоку вектора
$\phi\,\delta\Dfield$ крізь нескінченно віддалену поверхню:
\begin{equation*}
	\iiint\limits_V \divg\left(\phi\, \delta\Dfield\right) dV = \oiint\limits_{S\to\infty} \phi\, \delta\Dfield\cdot d\vect S.
\end{equation*}
\begin{Attention}
	Для будь-якої обмеженої в просторі системи зарядів на великих відстанях $\phi \sim 1/r$ і $D \sim 1/r^2$, тоді як площа сфери зростає як $r^2$.
	Тому підінтегральний вираз спадає як $1/r^3\cdot r^2 = 1/r \to 0$ при $r\to\infty$, і потік через нескінченно віддалену поверхню
	дорівнює нулю. Отже, перший доданок зникає тотожно.
\end{Attention}
У другому доданку скористаємось зв'язком напруженості поля з потенціалом, $\Efield = -\vect\nabla\phi$, і остаточно отримуємо
\begin{equation*}
	\delta W = \iiint\limits_V \frac{\Efield\cdot\delta\Dfield}{8\pi}\, dV.
\end{equation*}
Для лінійного діелектрика ($\Dfield = \epsilon\Efield$) це співвідношення інтегрується елементарно ($\Efield\cdot\delta\Dfield =
\tfrac12\delta(\Efield\cdot\Dfield)$), і ми приходимо до загального результату:
\begin{equation}\label{eq:GeneralEnergyDensity}
	\tcbhighmath{
		W = \iiint\limits_V \frac{\Efield\cdot\Dfield}{8\pi}\, dV, \qquad w = \frac{\Efield\cdot\Dfield}{8\pi} = \frac{\epsilon E^2}{8\pi}.
	}
\end{equation}

\begin{Attention}
	На відміну від виведення в попередньому пункті, отриманого лише для однорідного поля плоского конденсатора, формула~\eqref{eq:GeneralEnergyDensity}
	абсолютно загальна: вона справедлива для довільного, як завгодно неоднорідного розподілу заряду і довільної форми діелектричного середовища.
	Енергія поля --- це інтеграл від густини енергії $w = \Efield\cdot\Dfield/8\pi$ за всім об'ємом, де поле відмінне від нуля.
\end{Attention}

\subsection*{Власна та взаємна енергія зарядів}

Порівняємо два, на перший погляд, різні вирази для енергії системи зарядів: \enquote{дискретну} формулу~\eqref{eq:PairEnergy},
$W = \tfrac12\sum q_iq_j/r_{ij}$, і \enquote{польову} формулу~\eqref{eq:GeneralEnergyDensity}, $W = \int E^2/(8\pi)\, dV$. Перша з них, узагалі кажучи, може
набувати як додатних, так і від'ємних значень (притягання різнойменних зарядів дає від'ємний внесок); друга ж, як інтеграл від невід'ємної величини
$E^2$, --- \emphx{завжди} невід'ємна. Це здається суперечністю, проте насправді обидві формули описують не зовсім те саме.

Розглянемо поле системи двох зарядів, $\Efield = \Efield_1 + \Efield_2$. Тоді
\begin{equation*}
	W = \int\limits_V \frac{E^2}{8\pi}\, dV = \int\limits_V \frac{E_1^2}{8\pi}\, dV + \int\limits_V \frac{E_2^2}{8\pi}\, dV + \int\limits_V
	\frac{\Efield_1\cdot\Efield_2}{8\pi}\, dV.
\end{equation*}
Перші два доданки --- це \emphx{власна енергія} кожного із зарядів окремо; вони завжди додатні,
\begin{equation*}
	\int\limits_V \frac{E_1^2}{8\pi}\, dV > 0, \qquad \int\limits_V \frac{E_2^2}{8\pi}\, dV > 0,
\end{equation*}
і не залежать від взаємного розташування зарядів. Третій доданок --- це \emphx{взаємна (інтерференційна) енергія}
\begin{equation}\label{eq:SelfMutualEnergy}
	W_\text{вз} = \int\limits_V \frac{\Efield_1\cdot\Efield_2}{8\pi}\, dV,
\end{equation}
яка може мати будь-який знак залежно від взаємного розташування й знаків зарядів, і саме вона точно відповідає \enquote{дискретній} формулі
$q_1q_2/r_{12}$ з попереднього пункту.

\begin{Attention}
	Формула~\eqref{eq:PairEnergy} для точкових зарядів описує \emphx{лише взаємну} енергію --- енергію, потрібну для зближення вже готових точкових
	зарядів з нескінченності. Вона свідомо \emphx{не враховує} власну енергію кожного точкового заряду окремо, оскільки поле точкового заряду
	$E \sim 1/r^2$ поблизу нього необмежено зростає, а його власна енергія $\int E^2/(8\pi)\, dV$ --- розбігається (звертається в нескінченність) при
	$r \to 0$. Саме тому \enquote{дискретна} енергія системи точкових зарядів може бути від'ємною --- вона є лише частиною (взаємною частиною) від
	справді повної, завжди додатної, енергії поля, яка додатково містить (нескінченно великі) власні енергії кожного точкового заряду. Для тіл
	скінченних розмірів (як у прикладах з кулею вище) власна енергія, звісно, залишається скінченною.
\end{Attention}

%% --------------------------------------------------------
\section{Пондеромоторні сили}
%% --------------------------------------------------------

\subsection*{Метод віртуальних переміщень}

Одним з найефективніших способів обчислення сил, що діють на заряджені провідники й діелектрики, є \emphx{метод віртуальних переміщень}. Його ідея
полягає в тому, щоб обчислити роботу сил поля $\delta A = F_x\, dx$ при уявному (\emphx{віртуальному}) нескінченно малому зміщенні тіла на $dx$ у
напрямку, що нас цікавить, після чого шукана сила знаходиться просто як
\begin{equation*}
	F_x = \frac{\delta A}{dx}.
\end{equation*}

\begin{Attention}
	Метод віртуальних переміщень --- універсальний: він не вимагає з'ясування фізичної природи сили (електричної чи механічної) і автоматично
	враховує \emphx{всі} силові взаємодії в системі, оскільки спирається лише на закон збереження енергії.
\end{Attention}

\subsection*{Сила при сталому заряді провідників}

Нехай на провідниках підтримується сталий заряд, $q = \const$, тобто в системі немає зовнішніх джерел енергії (батарей тощо). За таких умов уся
робота $\delta A$ внутрішніх сил системи при віртуальному переміщенні відбувається виключно за рахунок \emphx{зменшення} власної електричної
енергії системи $W$:
\begin{equation*}
	\delta A = -\left(dW\right)_{q=\const}.
\end{equation*}
Якщо шукана сила $F_x$ здійснює роботу $\delta A = F_x\, dx$ при зміщенні тіла на $dx$, то
\begin{equation}\label{eq:ForceConstQ}
	\tcbhighmath{
		F_x = -\left.\frac{dW}{dx}\right|_{q=\const}.
	}
\end{equation}

\begin{Attention}
	Формула~\eqref{eq:ForceConstQ} --- це суто математичний прийом: для її застосування не обов'язково, щоб заряди провідників у \emphx{реальному}
	процесі залишалися незмінними. Достатньо знайти приріст енергії $dW$, уявляючи, що заряд саме такий, яким він є в даний момент, і \emphx{умовно}
	вважаючи його фіксованим при обчисленні похідної.
\end{Attention}

\subsection*{Сила при сталій різниці потенціалів}

Нехай тепер, навпаки, на провідниках підтримуються сталими потенціали, $\Delta\phi = \const$ (наприклад, провідники приєднані до батареї, яка
підживлює заряди в міру переміщення тіл і виконує при цьому роботу). За законом збереження енергії, робота джерела $\delta A_\text{дж}$ витрачається
одночасно на виконання механічної роботи $F_x\, dx$ і на зміну енергії поля $dW$:
\begin{equation*}
	\delta A_\text{дж} = F_x\, dx + dW.
\end{equation*}
З іншого боку, робота джерела з перенесення заряду $\delta q$ через різницю потенціалів $\Delta\phi$ дорівнює $\delta A_\text{дж} = \delta q\,
\Delta\phi$; оскільки $q = C\Delta\phi$, а $\Delta\phi = \const$, маємо $\delta q = \Delta\phi\, \delta C$, тобто $\delta A_\text{дж} = \Delta\phi^2\,
\delta C$. Водночас, за формулою~\eqref{eq:CapacitorEnergy}, $dW = \tfrac12\Delta\phi^2\, \delta C$ (енергія при сталому $\Delta\phi$ зростає разом
із ємністю). Підставляючи, знаходимо
\begin{equation*}
	F_x\, dx = \Delta\phi^2\,\delta C - \frac12\Delta\phi^2\,\delta C = \frac12\Delta\phi^2\,\delta C = dW,
\end{equation*}
звідки
\begin{equation}\label{eq:ForceConstPhi}
	\tcbhighmath{
		F_x = +\left.\frac{dW}{dx}\right|_{\Delta\phi=\const}.
	}
\end{equation}

\begin{Attention}
	Формули~\eqref{eq:ForceConstQ} і~\eqref{eq:ForceConstPhi} відрізняються \emphx{знаком}, проте описують \emphx{одну й ту саму} фізичну силу: сама
	сила, що діє на нерухоме (у момент обчислення) тіло, не може залежати від того, який уявний спосіб зміни системи ми обираємо для її обчислення
	--- метод сталого заряду чи метод сталого потенціалу. Різниця знаків пояснюється тим, що при сталому $\Delta\phi$ енергію в систему додатково
	вкладає джерело, і саме цей додатковий внесок \enquote{перевертає} знак похідної.
\end{Attention}

\subsection*{Приклад: сила тяжіння обкладок зарядженого конденсатора}

Знайдемо силу взаємного притягання обкладок плоского конденсатора, уявивши, що одна з обкладок зазнає віртуального зміщення $dx$ вздовж напрямку,
що з'єднує обкладки (рис.~\ref{pic:VirtualDisplacement}).

\begin{figure}[h!]\centering
	\localinput{virtual.tikz}
	\caption{Віртуальне зміщення $dx$ обкладки конденсатора\label{pic:VirtualDisplacement}}
\end{figure}

\emphx{Спосіб 1: $q = \const$.} Зміна енергії поля при зміні відстані між обкладками:
\begin{equation*}
	dW = d\left(\frac{q^2}{2C}\right) = -\frac{q^2}{2C^2}\, dC, \qquad C = \frac{\epsilon S}{4\pi x}, \qquad dC = -\frac{C}{x}\, dx,
\end{equation*}
звідки $dW = \dfrac{q^2}{2Cx}\, dx$. Виражаючи через напруженість поля ($q = CEd = C E x$, де $x$ тут відіграє роль відстані між обкладками),
\begin{equation*}
	dW = \frac{E^2 x^2\, \epsilon S}{2 \cdot 4\pi x \cdot x}\, dx = \frac{\epsilon E^2 S}{8\pi}\, dx,
\end{equation*}
і за формулою~\eqref{eq:ForceConstQ}
\begin{equation}\label{eq:ForceCapacitorPlates}
	\tcbhighmath{
		F_x = -\frac{dW}{dx} = -\frac{\epsilon E^2}{8\pi}\, S.
	}
\end{equation}

\emphx{Спосіб 2: $\Delta\phi = \const$.} Аналогічно, $dW = d\left(\dfrac{C\Delta\phi^2}{2}\right) = \dfrac{\Delta\phi^2}{2}\, dC =
-\dfrac{\Delta\phi^2 C}{2x}\, dx$, що після підстановки $\Delta\phi = Ex$ дає $dW = -\dfrac{\epsilon E^2 S}{8\pi}\, dx$, і за
формулою~\eqref{eq:ForceConstPhi} знову отримуємо той самий результат~\eqref{eq:ForceCapacitorPlates}.

\begin{Attention}
	Знак \enquote{мінус} у формулі~\eqref{eq:ForceCapacitorPlates} свідчить про те, що сила спрямована на \emphx{зменшення} відстані між обкладками, тобто
	обкладки \emphx{притягуються} одна до одної --- незалежно від знаків зарядів на них, оскільки сила пропорційна $E^2$.
\end{Attention}

\subsection*{Приклад: сила втягування діелектрика в конденсатор}

Розглянемо плоский конденсатор, у який частково (на глибину $x$) всунуто пластину діелектрика з проникністю $\epsilon$, і знайдемо силу, що діє на
цю пластину. Зручно замінити реальну систему еквівалентною --- двома паралельно з'єднаними конденсаторами: один, площею $S_1$, --- вакуумний, а
інший, площею $S_2$, --- повністю заповнений діелектриком (рис.~\ref{pic:DielectricInCapacitor}).

\begin{figure}[h!]\centering
	\localinput{dielectricInCapacitor.tikz}
	\caption{Модель конденсатора з частково всунутою пластиною діелектрика: паралельне з'єднання
		вакуумного та заповненого діелектриком конденсаторів\label{pic:DielectricInCapacitor}}
\end{figure}

За формулою для паралельного з'єднання~\eqref{eq:ParallelCap},
\begin{equation*}
	C = C_1 + C_2 = \frac{S_1}{4\pi d} + \frac{\epsilon S_2}{4\pi d}.
\end{equation*}
Нехай повна довжина конденсатора (у напрямку висування) дорівнює $L$, ширина обкладок --- $h$, тоді повна площа $S = Lh$ стала, а
$S_1 = \left(1 - \dfrac{x}{L}\right) S$, $S_2 = \dfrac{x}{L}S$ (тут $x$ --- глибина занурення пластини діелектрика). Підставляючи, знаходимо ємність
як функцію $x$:
\begin{equation*}
	C(x) = \frac{S}{4\pi d}\left(1 + (\epsilon - 1)\frac{x}{L}\right).
\end{equation*}
Обчислимо силу при сталому заряді $q = \const$ за формулою~\eqref{eq:ForceConstQ}:
\begin{equation*}
	F_x = -\frac{dW}{dx} = -\frac{dW}{dC}\frac{dC}{dx}, \qquad \frac{dW}{dC} = -\frac{q^2}{2C^2}, \qquad \frac{dC}{dx} = \frac{S}{4\pi d}\,
	\frac{\epsilon-1}{L},
\end{equation*}
звідки
\begin{equation*}
	F_x = \frac{q^2}{2C^2}\cdot\frac{S}{4\pi d}\cdot\frac{\epsilon - 1}{L} \stackrel{q = CEd}{=} \frac{(\epsilon-1)E^2}{8\pi}\, hd,
\end{equation*}
де ми скористались тим, що в кожен момент $q = CEd$, а $S/L = h$. Остаточно,
\begin{equation}\label{eq:ForceDielectric}
	\tcbhighmath{
		F_x = \frac{(\epsilon-1)E^2}{8\pi}\, S_\text{торця},
	}
\end{equation}
де $S_\text{торця} = hd$ --- площа торцевого перерізу пластини діелектрика (перерізу, яким вона \enquote{входить} у конденсатор). Оскільки $\epsilon > 1$,
сила $F_x$ додатна --- вона спрямована в бік \emphx{втягування} діелектрика всередину конденсатора.

\begin{Attention}
	Результат~\eqref{eq:ForceDielectric} має простий фізичний зміст: діелектрик завжди \emphx{втягується} в область сильнішого поля, оскільки
	заповнення проміжку поляризовним середовищем енергетично вигідне --- воно збільшує ємність (при сталому заряді енергія при цьому \emphx{падає}),
	а система, як завжди в природі, \enquote{прагне} до стану з меншою енергією. Це той самий фізичний механізм, завдяки якому рідкий діелектрик
	підіймається між частково зануреними в нього обкладками зарядженого конденсатора.
\end{Attention}

%% --------------------------------------------------------
\section{Підсумки розділу}
%% --------------------------------------------------------

\subsection*{Електрична ємність}
\begin{itemize}
	\item Ємність провідника: $\displaystyle C = \frac{q}{\phi}$; взаємна ємність пари провідників: $\displaystyle C = \frac{q}{\Delta\phi}$.
	\item Послідовне з'єднання конденсаторів: $\displaystyle \frac1C = \sum\limits_{i=1}^n \frac1{C_i}$; паралельне: $\displaystyle C =
	      \sum\limits_{i=1}^n C_i$.
	\item Ємність усамітненої кулі: $C = R$; плоского конденсатора: $\displaystyle C = \frac{\epsilon S}{4\pi d}$; сферичного: $\displaystyle C =
	      \frac{\epsilon R_1R_2}{R_2-R_1}$; циліндричного: $\displaystyle C = \frac{\epsilon\ell}{2\ln(R_2/R_1)}$.
\end{itemize}

\subsection*{Енергія електричного поля}
\begin{itemize}
	\item Взаємна енергія системи точкових зарядів: $\displaystyle U = \frac12 \sum_{i}\sum_{j\neq i} \frac{q_iq_j}{r_{ij}} = \frac12\sum_i
	      q_i\phi_i$.
	\item Енергія тіла з неперервним розподілом заряду: $\displaystyle W = \frac12\iiint\limits_V \rho\phi\, dV + \frac12\iint\limits_S \sigma\phi\,
	      dS$.
	\item Енергія зарядженого конденсатора: $\displaystyle W = \frac{q^2}{2C} = \frac{C\Delta\phi^2}{2}$.
	\item Об'ємна густина енергії поля (загальна формула): $\displaystyle w = \frac{\Efield\cdot\Dfield}{8\pi} = \frac{\epsilon E^2}{8\pi}$.
\end{itemize}

\subsection*{Пондеромоторні сили}
\begin{itemize}
	\item Метод віртуальних переміщень: $\displaystyle F_x = -\left.\frac{dW}{dx}\right|_{q=\const} = +\left.\frac{dW}{dx}\right|_{\Delta\phi=\const}$.
	\item Сила притягання обкладок конденсатора: $\displaystyle F_x = -\frac{\epsilon E^2}{8\pi}\, S$.
	\item Сила втягування діелектрика в конденсатор: $\displaystyle F_x = \frac{(\epsilon-1)E^2}{8\pi}\, S_\text{торця}$.
\end{itemize}

\begin{Attention}
	Усі результати цього розділу спираються на одну спільну ідею: електричне поле --- не допоміжна розрахункова абстракція, а фізичний об'єкт, який
	\emphx{сам є носієм енергії}. Ємність описує, скільки заряду (а отже, і поля) система здатна накопичити при заданій різниці потенціалів; енергія
	поля визначається інтегралом густини $w = \Efield\cdot\Dfield/8\pi$ за областю, де це поле існує; а пондеромоторні сили, що діють на провідники й
	діелектрики, --- це прояв фундаментального прагнення будь-якої фізичної системи до мінімуму енергії.
\end{Attention}
